\documentclass[a4paper]{arrowhead}

\usepackage[yyyymmdd]{datetime}
\usepackage{etoolbox}
\usepackage[T1]{fontenc}
\usepackage[utf8]{inputenc}
\usepackage{multirow}
\usepackage{longtable}
\usepackage{amssymb}
\usepackage{float}
\usepackage{tikz}
\usetikzlibrary{positioning,arrows.meta,fit,calc,shapes.misc,shapes.geometric}

\renewcommand{\dateseparator}{-}
\setlength{\emergencystretch}{3em}
\lstset{basicstyle=\ttfamily\footnotesize,breaklines=true,columns=fullflexible,frame=single,rulecolor=\color{gray!50},backgroundcolor=\color{gray!6}}
\tikzset{
  ahsys/.style={draw=ArrowheadBlue,very thick,fill=ArrowheadBlue!6,rounded corners=2pt,minimum width=3.2cm,minimum height=1.1cm,align=center,font=\small\bfseries},
  ahext/.style={draw=ArrowheadGrey,thick,dashed,fill=gray!5,rounded corners=2pt,minimum width=3.2cm,minimum height=1.0cm,align=center,font=\small},
  ahsrv/.style={draw=ArrowheadPurple,thick,fill=ArrowheadPurple!8,rounded corners=6pt,minimum height=0.55cm,inner xsep=4pt,align=center,font=\scriptsize\ttfamily},
  ahcsm/.style={draw=ArrowheadTeal,thick,fill=ArrowheadTeal!8,rounded corners=6pt,minimum height=0.55cm,inner xsep=4pt,align=center,font=\scriptsize\ttfamily},
  ahdb/.style={draw=ArrowheadOrange,thick,fill=ArrowheadOrange!8,cylinder,shape border rotate=90,aspect=0.25,minimum height=0.9cm,align=center,font=\scriptsize},
  ahflow/.style={-{Stealth[length=2.2mm]},thick,ArrowheadGrey},
  ahlabel/.style={font=\scriptsize,fill=white,inner sep=1pt,align=center},
}

%% Special references
\newcommand{\fref}[1]{{\textcolor{ArrowheadBlue}{\hyperref[sec:functions:#1]{#1}}}}
\newcommand{\mref}[1]{{\textcolor{ArrowheadPurple}{\hyperref[sec:model:#1]{#1}}}}
\newcommand{\sref}[1]{{\textcolor{ArrowheadPurple}{#1}}}
\newrobustcmd\msubsection[2]{
  \addtocounter{subsection}{1}
  \addcontentsline{toc}{subsection}{\protect\numberline{\thesubsection}#1 \textcolor{ArrowheadPurple}{#2}}
  \subsection*{\thesubsection\quad#1 \textcolor{ArrowheadPurple}{#2}}
  \label{sec:model:#2}
}

\begin{document}

%% Arrowhead Document Properties
\ArrowheadTitle{WorkflowConfigurator}

\ArrowheadType{System Design Description}
\ArrowheadTypeShort{SysDD}
\ArrowheadVersion{5.4.0}
\ArrowheadDate{2026-11-24}
\ArrowheadAuthor{Documentation maintenance}
\ArrowheadStatus{REVIEW DRAFT}
\ArrowheadContact{Repository maintainers}
\ArrowheadFooter{\href{www.arrowhead.eu}{www.arrowhead.eu}}
\ArrowheadSetup
%%

%% Front Page
\begin{center}
  \vspace*{1cm}
  \huge{\arrowtitle}

  \vspace*{0.2cm}
  \LARGE{\arrowtype}
  \vspace*{1cm}

  
  \vspace*{\fill}

  

  \vspace*{1cm}
  \vspace*{\fill}

  % Front Page Abstract
  \begin{abstract}
    This document describes how the WorkflowConfigurator is implemented in the Java reference implementation: its internal components, the realisation of its functional and non-functional properties, its configuration and the implementation status of the other languages. Documentation edition 5.4.0.
  \end{abstract}

  \vspace*{1cm}
\end{center}

\newpage
%%

%% Table of Contents
\tableofcontents
\newpage
%%

\section{Overview}
\label{sec:overview}
This document describes the design of the WorkflowConfigurator, which lets an engineer draw workflows with conditions between the microservices of configured local clouds, and validates and stores them in the WorkflowManager. The abstract system is described in the
SysD of the same name; its produced services are specified by SDs and IDDs in this documentation set.

\begin{figure}[ht!]
  \centering
  \resizebox{\textwidth}{!}{
\begin{tikzpicture}[node distance=4mm]
\node[ahsys,minimum height=6.00cm,minimum width=3.6cm] (sys) at (0,0) {WorkflowConfigurator\\[-1pt]{\normalfont\scriptsize $\ll$system$\gg$}};
\node[ahsrv,anchor=west] (p0) at (3.4,2.62) {workflowConfiguration};
\draw[thick,ArrowheadPurple] (1.8,2.62) -- (p0.west);
\fill[ArrowheadPurple] (1.8,2.62) circle (1.4pt);
\node[ahsrv,anchor=west] (p1) at (3.4,1.88) {generalManagement};
\draw[thick,ArrowheadPurple] (1.8,1.88) -- (p1.west);
\fill[ArrowheadPurple] (1.8,1.88) circle (1.4pt);
\node[ahcsm,anchor=east] (c0) at (-3.4,2.62) {serviceDiscovery\\[-2pt]{\tiny\rmfamily ServiceRegistry}};
\draw[thick,ArrowheadTeal,-{Stealth[length=1.6mm]}] (c0.east) -- (-1.8,2.62);
\node[ahcsm,anchor=east] (c1) at (-3.4,1.88) {serviceRegistryManagement\\[-2pt]{\tiny\rmfamily ServiceRegistry}};
\draw[thick,ArrowheadTeal,-{Stealth[length=1.6mm]}] (c1.east) -- (-1.8,1.88);
\node[ahcsm,anchor=east] (c2) at (-3.4,1.12) {identity\\[-2pt]{\tiny\rmfamily Authentication}};
\draw[thick,ArrowheadTeal,-{Stealth[length=1.6mm]}] (c2.east) -- (-1.8,1.12);
\node[ahcsm,anchor=east] (c3) at (-3.4,0.38) {authorization\\[-2pt]{\tiny\rmfamily ConsumerAuthorization}};
\draw[thick,ArrowheadTeal,-{Stealth[length=1.6mm]}] (c3.east) -- (-1.8,0.38);
\node[ahcsm,anchor=east] (c4) at (-3.4,-0.38) {workflowRecipe\\[-2pt]{\tiny\rmfamily WorkflowManager}};
\draw[thick,ArrowheadTeal,-{Stealth[length=1.6mm]}] (c4.east) -- (-1.8,-0.38);
\node[ahcsm,anchor=east] (c5) at (-3.4,-1.12) {localCloudConfiguration\\[-2pt]{\tiny\rmfamily LocalCloudConfigurator}};
\draw[thick,ArrowheadTeal,-{Stealth[length=1.6mm]}] (c5.east) -- (-1.8,-1.12);
\node[ahcsm,anchor=east] (c6) at (-3.4,-1.88) {soLocalCloudConfiguration\\[-2pt]{\tiny\rmfamily SoLocalCloudConfigurator}};
\draw[thick,ArrowheadTeal,-{Stealth[length=1.6mm]}] (c6.east) -- (-1.8,-1.88);
\node[ahcsm,anchor=east] (c7) at (-3.4,-2.62) {sysMonitor (AMQP)\\[-2pt]{\tiny\rmfamily SysMonitor broker}};
\draw[thick,ArrowheadTeal,-{Stealth[length=1.6mm]}] (c7.east) -- (-1.8,-2.62);
\node[font=\scriptsize\bfseries,ArrowheadPurple] at (5.0,3.52) {produced};
\node[font=\scriptsize\bfseries,ArrowheadTeal] at (-5.0,3.52) {consumed};
\end{tikzpicture}
  }
  \caption{Block view of the WorkflowConfigurator with its produced and consumed services.}
  \label{fig:block}
\end{figure}


\newpage
\section{Implementation}
\label{sec:implementation}

\subsection{Implementation language and tools}
Java~21 and Spring Boot~4 (Spring Framework~7, Spring Security~7), Spring Data JPA with Hibernate, Jackson for JSON,
Springdoc for the OpenAPI description, built with Maven~3.9 (in the \texttt{maven:3.\allowbreak{}9-\allowbreak{}eclipse-\allowbreak{}temurin-\allowbreak{}21} image). The
system builds on the vendored Arrowhead~5.4.0 common library (\texttt{vendor/\allowbreak{}common-\allowbreak{}utils},
\texttt{vendor/\allowbreak{}data-\allowbreak{}transfer-\allowbreak{}objects}): system registration, identity, TLS, management filters and HTTP consumption
of other services. The source is module \texttt{workflow-\allowbreak{}configurator} of \texttt{Eclipse-\allowbreak{}Arrowhead-\allowbreak{}java/\allowbreak{}ah5-\allowbreak{}Workflow-\allowbreak{}java}; the
container image runs the Spring Boot jar on Eclipse Temurin~21.

\subsection{Internal structure}

{\small
\begin{longtable}{|p{6.2cm}|p{9.6cm}|}

\hline
\rowcolor{gray!33} Component & Responsibility \\ \hline
\endhead
{\footnotesize \texttt{Workflow\allowbreak{}Configurator\allowbreak{}System\allowbreak{}Info}} & Registration of \texttt{workflow\allowbreak{}Configuration} and \texttt{general\allowbreak{}Management}. \\ \hline
{\footnotesize \texttt{api.\allowbreak{}http.\allowbreak{}ConfigurationAPI}} & The seven operations; upstream failures returned with the provider and status. \\ \hline
{\footnotesize \texttt{service.\allowbreak{}CloudDirectory}} & Clouds, systems, services and executors from the ServiceRegistry and the configurators. \\ \hline
{\footnotesize \texttt{service.\allowbreak{}RecipeGateway,\allowbreak{} Upstream}} & Calls to the WorkflowManager and other providers through the common library. \\ \hline
{\footnotesize \texttt{console.\allowbreak{}ConsoleConnectorConfig,\allowbreak{} ConsoleFilter,\allowbreak{} ConsoleAuth,\allowbreak{} ConsolePasswordHash,\allowbreak{} ConsoleController}} & Second Tomcat connector for the console; password login (scrypt), \texttt{wfc\_session} cookie, Host/Origin checks, lockout after five failures; static editor pages only on this connector. \\ \hline
{\footnotesize \texttt{security.\allowbreak{}ConfiguratorSecurityConfig,\allowbreak{} ServiceAuthorizationFilter,\allowbreak{} ConfiguratorOperations}} & HTTPS redirect limited to the Arrowhead connector; provider-side authorization of the seven operations. \\ \hline
{\footnotesize \texttt{resources/\allowbreak{}web: index.\allowbreak{}html,\allowbreak{} app.\allowbreak{}js,\allowbreak{} model.\allowbreak{}js,\allowbreak{} app.\allowbreak{}css}} & The editor: rendering and interaction (app.js), recipe model, local checks, layout and help (model.js), tested with Node tests. \\ \hline
\end{longtable}
}


\subsection{Functional properties implementation}
The console connector serves the editor and the same \texttt{/\allowbreak{}workflowconfigurator/\allowbreak{}api} operations to an
authenticated browser session; requests from the console are marked as operator requests and run the service logic
without an mTLS caller, while the Arrowhead connector returns 404 for console pages. The editor's model module mirrors
the WorkflowManager's checks for immediate feedback and uses the same codes, so every finding links to the same help
text; the WorkflowManager's verdict remains authoritative and saving requires ACCEPT. Arrange computes column gaps
from the widest condition label and places labels on their transition before the target step, moving them with a
leader line when the spot is taken.

\subsection{Non functional properties implementation}
\subsubsection{Security}
Spring Security filter chains of the common library: TLS client authentication, identity-token verification against
the Authentication system (outsourced authentication), the management filter and the workflow systems'
\texttt{Service\allowbreak{}Authorization\allowbreak{}Filter}, which maps method and path to service definition and operation and asks
ConsumerAuthorization (\texttt{authorization.\allowbreak{}verify}). Object-level checks (holder, owner, configured manager) are made in
the service logic and give 403.
\subsubsection{Power management}
No specific power management. Thread pools are bounded; the WorkflowManager's watchdog and the executor's pull loop are
the only periodic activities.
\subsubsection{Internal monitoring}
Spring Boot health endpoint,
the common-library log table exposed through \texttt{general\allowbreak{}Management.\allowbreak{}get-\allowbreak{}log}, and optional sysMonitor telemetry over
AMQP (\texttt{sysmonitor.\allowbreak{}enabled}).
\subsubsection{Configuration}

{\small
\begin{longtable}{|p{5.6cm}|p{3.4cm}|p{6.8cm}|}

\hline
\rowcolor{gray!33} Property & Default & Meaning \\ \hline
\endhead
\texttt{server.\allowbreak{}port} & 8484 & Arrowhead connector port. \\ \hline
\texttt{console.\allowbreak{}port} & 8485 & Console connector port. \\ \hline
\texttt{console.\allowbreak{}address} & 127.0.0.1 & Console bind address (the container deployment binds 0.0.0.0 and publishes to host loopback). \\ \hline
\texttt{console.\allowbreak{}password-\allowbreak{}hash} & (secret) & \texttt{scrypt:<salt>:<hash>} (N=16384, r=8, p=1); the console is closed without it. \\ \hline
\texttt{console.\allowbreak{}allowed-\allowbreak{}hosts} & localhost,127.0.0.1,[::1] & Accepted \texttt{Host} headers. \\ \hline
\texttt{console.\allowbreak{}session-\allowbreak{}seconds} & 28800 & Console session lifetime. \\ \hline
\texttt{workflow.\allowbreak{}configurator.\allowbreak{}manager} & WorkflowManager & Recipe store. \\ \hline
\texttt{workflow.\allowbreak{}configurator.\allowbreak{}local-\allowbreak{}cloud-\allowbreak{}configurator} & LocalCloudConfigurator & Design source. \\ \hline
\texttt{workflow.\allowbreak{}configurator.\allowbreak{}solc-\allowbreak{}configurator} & So\allowbreak{}Local\allowbreak{}Cloud\allowbreak{}Configurator & SoLC design source. \\ \hline
\texttt{workflow.\allowbreak{}configurator.\allowbreak{}solc.\allowbreak{}enabled} & true & List SoLC member clouds. \\ \hline
\texttt{workflow.\allowbreak{}configurator.\allowbreak{}hidden-\allowbreak{}services} & sys\allowbreak{}Monitor,\allowbreak{}general\allowbreak{}Management & Services not offered in the palette. \\ \hline
\texttt{workflow.\allowbreak{}standalone} & false & Development mode without core systems. \\ \hline
\texttt{service.\allowbreak{}authorization.\allowbreak{}enabled} & false (deployment: true) & Provider-side ConsumerAuthorization check. \\ \hline
\end{longtable}
}


\newpage
\section{Services}
\label{sec:services}
\subsection{Produced services}
All produced services use interface template \texttt{generic\_\allowbreak{}https}, policy \texttt{CERT\_AUTH} and version 1.0.0 and are
registered with every operation's method and path.

{\small
\begin{longtable}{|p{3.8cm}|p{6.2cm}|p{5.8cm}|}

\hline
\rowcolor{gray!33} Service & Operations & Documents \\ \hline
\endhead
\texttt{workflow\allowbreak{}Configuration} & list-clouds, list-services, list-executors, list-recipes, get-recipe, validate-recipe, store-recipe & \texttt{Workflow\allowbreak{}Configurator-\allowbreak{}v5.\allowbreak{}4.\allowbreak{}0-\allowbreak{}workflow\allowbreak{}Configuration-\allowbreak{}SD}, \texttt{-HTTP-IDD} \\ \hline
\texttt{general\allowbreak{}Management} & get-log, get-config & \texttt{Workflow-\allowbreak{}v5.\allowbreak{}4.\allowbreak{}0-\allowbreak{}general\allowbreak{}Management-\allowbreak{}SD}, \texttt{-HTTP-IDD} \\ \hline
\end{longtable}
}

\subsection{Consumed services}

{\small
\begin{longtable}{|p{3.4cm}|p{3.2cm}|p{3.6cm}|p{5.6cm}|}

\hline
\rowcolor{gray!33} Service & Provider & Operations & Purpose \\ \hline
\endhead
\texttt{service\allowbreak{}Discovery} & ServiceRegistry & register, revoke, lookup & Registration; executor and fallback discovery. \\ \hline
\texttt{service\allowbreak{}Registry\allowbreak{}Management} & ServiceRegistry & system-query, service-query & Systems and services of the live cloud (needs a grant and the \texttt{authorization} management policy). \\ \hline
\texttt{identity} & Authentication & login, logout, identity-verify & Own identity; verification of callers. \\ \hline
\texttt{authorization} & ConsumerAuthorization & verify & Provider-side authorization. \\ \hline
\texttt{workflow\allowbreak{}Recipe} & WorkflowManager & validate, register, query, get & Validation and storage. \\ \hline
\texttt{local\allowbreak{}Cloud\allowbreak{}Configuration} & LocalCloudConfigurator & catalog, list-configurations & Local-cloud designs. \\ \hline
\texttt{so\allowbreak{}Local\allowbreak{}Cloud\allowbreak{}Configuration} & So\allowbreak{}Local\allowbreak{}Cloud\allowbreak{}Configurator & list-designs, get-design & SoLC member clouds (when \texttt{workflow.\allowbreak{}configurator.\allowbreak{}solc.\allowbreak{}enabled}). \\ \hline
\texttt{sysMonitor (AMQP)} & SysMonitor broker & publish & When \texttt{sysmonitor.\allowbreak{}enabled}. \\ \hline
\end{longtable}
}


\section{Verification}
\label{sec:verification}
Seven Spring tests and 12 Node tests of the editor model (\texttt{node --test}
\texttt{workflow-\allowbreak{}configurator/\allowbreak{}src/\allowbreak{}test/\allowbreak{}js/\allowbreak{}model.\allowbreak{}test.\allowbreak{}mjs}), all passed on 2026-11-24, among them a test that every validator code has help text. In the
cloud test the configurator lists the running cloud, three LocalCloudConfigurator designs and five SoLC member clouds,
refuses a condition reading a later step (422, G03), stores a guarded recipe (201, then 200 for identical content and
409 for different content) and enforces ConsumerAuthorization at the Arrowhead connector. The editor was also driven in
a browser to draw, validate, save and reopen a branching workflow, which the WorkflowManager then ran correctly.


\section{Implementation variants}
\label{sec:variants}
\begin{itemize}
\item \textbf{Java (reference implementation).} \texttt{Eclipse-\allowbreak{}Arrowhead-\allowbreak{}java/\allowbreak{}ah5-\allowbreak{}Workflow-\allowbreak{}java/\allowbreak{}workflow-\allowbreak{}configurator}, Java 21, Spring
Boot, built with Maven together with the shared \texttt{workflow-\allowbreak{}model} module and the vendored Arrowhead~5.4.0 common
library.
\item \textbf{Go.} Not implemented. No workflow system exists in \texttt{Eclipse-\allowbreak{}Arrowhead-\allowbreak{}go}.
\item \textbf{Rust.} Not implemented. No workflow system exists in \texttt{Eclipse-\allowbreak{}Arrowhead-\allowbreak{}rust}.
\end{itemize}
The workflow systems are therefore not yet interoperable across the three implementations. A Go or Rust port must
reproduce the IDDs of this set exactly and, for the WorkflowManager, the behaviour of \texttt{workflow-\allowbreak{}model}: the
recipe schema and decoder, the validator's findings and verdicts, the guard grammar and the transition rules. The
architecture prescribes conformance vectors (recipe in, verdict and transition trace out) to be run by all three
implementations; they are not written yet.


\newpage
\bibliographystyle{IEEEtran}
\bibliography{bibliography}

\newpage

\section{Revision History}
\subsection{Amendments}

\noindent\begin{tabularx}{\textwidth}{| p{1cm} | p{2.2cm} | p{1.6cm} | X | p{3.6cm} |} \hline
\rowcolor{gray!33} No. & Date & Version & Subject of Amendments & Author \\ \hline
1 & 2026-11-23 & \arrowversion & Initial edition: system design description of the WorkflowConfigurator, written from the Java reference implementation (ah5-Workflow-java). & Documentation maintenance \\ \hline
2 & 2026-11-24 & \arrowversion & Verification status. & Documentation maintenance \\ \hline
\end{tabularx}

\subsection{Quality Assurance}

\noindent\begin{tabularx}{\textwidth}{| p{1cm} | p{2.2cm} | p{1.6cm} | X |} \hline
\rowcolor{gray!33} No. & Date & Version & Approved by \\ \hline
1 & & \arrowversion & Pending review \\ \hline
\end{tabularx}

\end{document}
